\documentclass{jsarticle}
\usepackage{amsmath}
\usepackage{amssymb}
\begin{document}
\title{}
\author{}
\date{}
\maketitle




作為体験  させられ体験、過去の妄想からの誘導と幻覚で精神が
特定の人から誘導されて生じていると想っている症状、その誘導されている
と思っているから、作為体験として実体験する妄想、
体験、思考の流れがさせられ体験として誘導されていると
妄想として体験している実体験\\
妄想追想　過去の妄想をあとから追っかけて関係づけてしまう妄想\\ 
放送伝播　自分の思考が他者に伝わっていると想う原因を放送されていると想う妄想\\ 
考察伝播　考えているから他者に思考が伝わっていると想っている妄想\\ 
広島大学附属病院にはいろいろな体験をしていた。\\
実体験とは現実ではあり得ない異常体験を実際に体験する妄想\\ 
能力で見ている人と、病気で見られていると想っている症状がある。\\ 
速読法はこの症状の緩和のいい援助にもなっている。\\ 
瞑想法も先生がすごくいいよと勧めるのも分かる。 \\ 
レボトミンは全脳瞑想ですごく楽だった。何も考えが生じてこなかった。\\
リスペリドンの液剤は、いくら飲んでも楽だった。\\
他者に思考が読まれている幻覚がなくなっていた。\\ 
補聴器が動いているのとは、幻覚の触手ないので、疑心暗鬼になっていただけだった。\\ 
リスペリドンは、幻覚、妄想、補聴器が思考から消えていた。\\ 
薬を飲み過ぎてはいけないのかなと想っていただけだった。\\

\end{document}
